\documentclass[11pt]{article}
\usepackage{amsmath,amssymb}
\usepackage[hidelinks]{hyperref}
\title{A scouting move changes both the observation and the time left to use it}
\author{}\date{}
\begin{document}\maketitle

\paragraph{Declared experiment.}
A scout and a rescue carrier in a small simulated warehouse must choose
one of two inspection bays. Exactly one bay, left or right, is occupied
by a stationary stopped agent during the cut. The scout starts at
\((2,3)\), the carrier below it at \((2,4)\). Both initial local crops
are identical under the two fault placements; placement and both
candidate goal cells are evaluator data, not given to the policy as an
answer. A separate carrier has a mandatory local move and uplink at tick 2.
The geometry, task rights, tick budget, stationary fault envelope and
links are authored; selected RWARE motion and its actual local image
extractor, plus selected SubT delivery and byte debit, execute the
effects. This is not a physical sensing model or a joint field trial.

\paragraph{Acquisition and access.}
On tick 1 the scout can pick up a scheduled optional shelf at its
starting cell or try moving one cell forward. The physical step to
\((2,2)\) changes the stock radius-one crop and reveals both candidate
cells. Only the occupancy layer of the scout's own cropped image
produces the short typed branch cue; the policy is not supplied the
actual faulty side. A failed move against another stopped occupant
does not create the observation. An erased crop also creates no cue.
If the send itself is erased, the scout pays the actual transmit debit
without producing the rescuer's receipt. If the cue is only private to
the scout, no early branch follows. The scout gives up the optional
pickup even when active acquisition or transmission fails. Different
failures are therefore neither a free observation nor a single
common observability state.

\paragraph{Motion and time of use.}
With the installed side-branch routes, carrier choices begin at tick 3.
A cue must reach that choice; the exact cue delivered at tick 4 is
already too late for those routes. But when the scout leaves \((2,3)\)
it also frees the carrier's first common forward move. The carrier
can verify the open cell from its own native local crop and move there
at tick 3. It may then use the same cue delivered at tick 4 and still
reach the required bay at tick 8. A single physical probe has altered
both the available evidence and the set of motions still possible
before commitment. Treating sensor latency and feasible common prefix
as independent would miss this recovery. If the shared prefix is
forbidden, it loses that extra tick again.

For a separate exact model check, enumerate all reachable oriented
states at the end of a shared prefix that cannot branch on the fault,
then for each one enumerate the admissible suffixes with a cue.
Require a suffix to the designated open bay in \emph{each} placement.
There are six unit action slots (ticks 3--8). When the common center
step is disallowed, zero slots before the cue is feasible and one or
more is not. If allowed, one uncontrolled-by-fault prefix slot is
possible; two or more is not. This finite optimality result is indexed
to the given motion rules and no extra observations beyond the cue
before commitment. It is not a bound against controllers with
additional free sensors, a different legal action set or prepositioning.

\paragraph{Matched comparison and service scope.}
Ninety-six selected native comparisons across six policies show the
early and flexible-late arms meeting the rescue deadline on both
placements in the qualified envelope (successful move, sampled crop,
and delivered cue). The late fixed route and private-only controls
fail that deadline, while unilateral left choice works in one
placement and not the other. Raw hit counts across the deliberately
failed tests are not empirical rates and do not rank expected utility.
All scheduled independent local jobs finish. The released side routes
and the common-prefix policy are ordinary active control. A direct
finite reachability construction gives the same frontier. Relevant
prior formulations include
\href{https://arxiv.org/abs/1410.0083}{active sensing for reactive synthesis}
and
\href{https://arxiv.org/abs/1209.1695}{partial-history sharing}.
No superiority over these planning principles is inferred. The result
adds an actual source-qualified travel/observation chain and distinguishes
late observation from lack of usable motion. Source of actuator release,
unmodeled optical effects, receiver authenticity, traffic under
fully shared physical routing, and independent mission acceptance remain
outside this envelope.
\end{document}
